% ============================================================
% Capítulo 10 — Entrenamiento: conceptos + código
% ============================================================
\chapter{Entrenamiento: optimizador, scheduler y mixed precision}
\label{cap:entrenamiento}

Este capítulo explica los conceptos detrás del loop de
entrenamiento y muestra el código línea por línea.

\section{¿Qué es AdamW?}

\subsection{Adam básico}

Adam (\emph{Adaptive Moment Estimation})~\cite{kingma2014adam}
es un optimizador que mantiene dos estadísticas por parámetro:

\begin{itemize}
  \item $m_t$: promedio exponencial del gradiente (momento 1).
  \item $v_t$: promedio exponencial del gradiente al cuadrado
        (momento 2).
\end{itemize}

\begin{align}
  m_t &= \beta_1 m_{t-1} + (1 - \beta_1) g_t \\
  v_t &= \beta_2 v_{t-1} + (1 - \beta_2) g_t^2 \\
  \hat{m}_t &= \frac{m_t}{1 - \beta_1^t} \\
  \hat{v}_t &= \frac{v_t}{1 - \beta_2^t} \\
  \theta_{t+1} &= \theta_t - \eta \cdot \frac{\hat{m}_t}{\sqrt{\hat{v}_t} + \epsilon}
\end{align}

\noindent\textbf{Ventaja}: adapta el learning rate por
parámetro, lo que acelera la convergencia en problemas con
gradientes de magnitudes muy distintas.

\subsection{AdamW}

AdamW~\cite{loshchilov2019adamw} separa el \emph{weight decay}
del optimizador, en vez de aplicarlo como una regularización L2
en la loss. Esto da una regularización más efectiva.

En este proyecto:
\begin{itemize}
  \item $\beta_1 = 0.90$ (\texttt{config.adam\_alpha}).
  \item $\beta_2 = 0.999$ (\texttt{config.adam\_beta}).
  \item lr = 0.001 (\texttt{config.lr}).
\end{itemize}

\section{¿Qué es un LR scheduler?}

Un \emph{learning rate scheduler} modifica el learning rate
durante el entrenamiento. Hay varios tipos:

\begin{itemize}
  \item \textbf{Step}: decae en epochs fijos.
  \item \textbf{Exponential}: decae exponencialmente.
  \item \textbf{Cosine}: cae siguiendo un coseno.
  \item \textbf{ReduceOnPlateau}: decae si la loss no mejora.
\end{itemize}

\noindent Este proyecto usa \texttt{MultiStepLR}: decae por un
factor de 0.1 en los epochs especificados en
\texttt{epoch\_decay = [20, 40]}.

\begin{itemize}
  \item Epochs 1--19: lr = 0.001.
  \item Epochs 20--39: lr = 0.0001.
  \item Epochs 40--60: lr = 0.00001.
\end{itemize}

\noindent\textbf{¿Por qué decaer?} Al principio se quiere
aprender rápido. Cerca del óptimo, lr grande hace oscilar;
reducirlo permite converger.

\section{¿Qué es mixed precision (FP16/BF16)?}

\subsection{Precisiones disponibles}

\begin{itemize}
  \item \textbf{FP32} (32 bits): precisión estándar. Ocupa 4
        bytes por número.
  \item \textbf{FP16} (16 bits): mitad de memoria. Acelera en
        GPUs con Tensor Cores (NVIDIA desde Volta).
  \item \textbf{BF16} (16 bits): mismo rango que FP32 pero
        menos precisión. \emph{No necesita scaler}.
\end{itemize}

\subsection{El problema del FP16}

FP16 tiene rango muy limitado ($\pm 65504$). Si un gradiente es
muy pequeño ($< 6 \cdot 10^{-5}$), se redondea a 0 y
\emph{desaparece}.

\subsection{Solución: GradScaler}

\texttt{torch.cuda.amp.GradScaler} multiplica la loss por un
factor (típicamente $2^{16}$) antes del \texttt{backward}, lo
que efectivamente \emph{amplifica} los gradientes. Luego los
des-escala al hacer \texttt{step}:

\begin{minted}[language=Python, basicstyle=\ttfamily\small, frame=single]{python}
scaler = GradScaler()
scaler.scale(loss).backward()      # loss * factor -> gradientes grandes
scaler.step(optimizer)             # des-escala y aplica
scaler.update()                    # ajusta el factor dinámicamente
\end{minted}

\subsection{Autocast}

\texttt{torch.autocast} aplica casting automático de tipos:

\begin{itemize}
  \item Operaciones \emph{seguras} (matmul, conv): en FP16.
  \item Operaciones \emph{sensibles} (log, softmax, loss): en
        FP32.
\end{itemize}

\noindent Esto se aplica con un context manager:

\begin{minted}[language=Python, basicstyle=\ttfamily\small, frame=single]{python}
with torch.autocast(device_type='cuda', dtype=torch.float16):
    output = model(input)
    loss = criterion(output, target)
\end{minted}

\subsection{GPU vs CPU}

En este proyecto:

\begin{minted}[language=Python, basicstyle=\ttfamily\small, frame=single]{python}
# GPU: FP16 con GradScaler
# CPU: BF16 sin scaler (más estable)
precision_dtype = torch.bfloat16 if config.device == 'cpu' \
                  else torch.float16
\end{minted}

\section{¿Qué hace \texttt{optimizer.zero\_grad()}?}

Por defecto PyTorch \emph{acumula} gradientes entre llamadas a
\texttt{backward()}. Hay que \textbf{limpiar} explícitamente los
gradientes al inicio de cada batch:

\begin{minted}[language=Python, basicstyle=\ttfamily\small, frame=single]{python}
optimizer.zero_grad()  # <- SIEMPRE antes de forward
output = model(input)
loss = criterion(output, target)
loss.backward()         # acumula en .grad
optimizer.step()        # aplica los gradientes
\end{minted}

\noindent\textbf{Alternativa}: \texttt{optimizer.zero\_grad(set\_to\_None=True)}
libera memoria al poner los gradientes en \texttt{None} en vez de
0.

\section{¿Cómo se guarda un checkpoint?}

Un checkpoint es un diccionario con todo lo necesario para
reanudar el entrenamiento:

\begin{minted}[language=Python, basicstyle=\ttfamily\small, frame=single]{python}
checkpoint = {
    'epoch': current_epoch,
    'model_state_dict': model.state_dict(),
    'optimizer_state_dict': optimizer.state_dict(),
    'loss': current_loss
}
torch.save(checkpoint, 'checkpoint.pth.tar')
\end{minted}

\noindent\textbf{state\_dict}: diccionario con todos los
parámetros del modelo (o del optimizer). Es lo que se mueve a
GPU/CPU con \texttt{model.load\_state\_dict()}.

\section{Código: \texttt{train.py} main + loop}

\begin{minted}[language=Python, numbers, firstnumber=34, linerange=34-49]{python}
def parse_args():
    parser = argparse.ArgumentParser(description='Train TextReIDNet on CUHK-PEDES')

    # Fuente del dataset: 'huggingface' (rápido) o 'original' (paper)
    parser.add_argument('--dataset_source', type=str, default='huggingface',
                        choices=['huggingface', 'original'])
    parser.add_argument('--epochs', type=int, default=None)
    parser.add_argument('--batch_size', type=int, default=None)
    parser.add_argument('--lr', type=float, default=None)
    parser.add_argument('--seed', type=int, default=None)
    parser.add_argument('--output_dir', type=str, default='./data/checkpoints')

    return parser.parse_args()
\end{minted}

\begin{minted}[language=Python, numbers, firstnumber=52, linerange=52-110]{python}
def main():
    args = parse_args()

    # Construir config desde config.py
    config = sys_configuration(
        dataset_name="CUHK-PEDES",
        dataset_source=args.dataset_source
    )

    # Aplicar overrides de CLI sobre la config
    if args.epochs is not None:
        config['epoch'] = args.epochs
    if args.batch_size is not None:
        config['batch_size'] = args.batch_size
    if args.lr is not None:
        config['lr'] = args.lr
    if args.seed is not None:
        config['seed'] = args.seed
    if args.output_dir is not None:
        config['model_save_path'] = os.path.abspath(args.output_dir)
        os.makedirs(config['model_save_path'], exist_ok=True)

    # Fijar semilla para reproducibilidad
    set_seed(config.seed)

    # Compartir memoria via filesystem (evita errores con muchos workers)
    torch.multiprocessing.set_sharing_strategy('file_system')

    # GradScaler para mixed precision (FP16/BF16)
    scaler = GradScaler()

    # Construir DataLoaders
    train_data_loader, _, _, train_num_classes = \
        build_cuhkpedes_dataloader(config)

    # Construir modelo y mover a device
    model = TextReIDNet(config).to(config.device)

    # IdentityLoss tiene parámetros entrenables (la capa Linear)
    identity_loss_fnx = IdentityLoss(
        config=config, class_num=train_num_classes
    ).to(config.device)

    # RankingLoss NO tiene parámetros entrenables
    ranking_loss_fnx = RankingLoss(config)

    # Optimizer: AdamW
    optimizer = optim.AdamW(
        model.parameters(),
        betas=(config.adam_alpha, config.adam_beta),
        lr=config.lr
    )

    # Scheduler: decae lr en epochs específicos
    scheduler = optim.lr_scheduler.MultiStepLR(optimizer, config.epoch_decay)
\end{minted}

\begin{minted}[language=Python, numbers, firstnumber=113, linerange=113-191]{python}
for current_epoch in range(1, config.epoch + 1):
    train_ranking_loss_list = []
    train_identity_loss_list = []
    train_total_loss_list = []

    model.train()

    with tqdm(train_data_loader, unit='batch') as tepoch:
        tepoch.set_description(f"Epoch {current_epoch}/{config.epoch}")

        for batch in tepoch:
            # Mover datos al device
            preprocessed_images = batch['preprocessed_images'].to(config.device)
            labels = batch['pids'].to(config.device)
            token_ids = batch['token_ids'].to(config.device)
            orig_token_length = batch['orig_token_lengths'].to(config.device)

            # Limpiar gradientes del batch anterior
            optimizer.zero_grad()

            # Mixed precision
            precision_dtype = torch.bfloat16 if config.device == 'cpu' \
                              else torch.float16
            with torch.autocast(device_type=config.device, dtype=precision_dtype):
                # FORWARD
                visual_embeddings, textual_embeddings = model(
                    image=preprocessed_images,
                    text_ids=token_ids,
                    text_length=orig_token_length
                )

                # Aplanar para losses
                visual_emb_flat = visual_embeddings.view(visual_embeddings.size(0), -1)
                text_emb_flat = textual_embeddings.view(textual_embeddings.size(0), -1)
                labels_flat = labels.view(-1)

                # Calcular losses
                ranking_loss = ranking_loss_fnx(visual_emb_flat, text_emb_flat, labels_flat)
                identity_loss = identity_loss_fnx(visual_embeddings, textual_embeddings, labels_flat)
                total_loss = (config.ranking_loss_alpha * ranking_loss
                              + config.identity_loss_beta * identity_loss)

            # BACKWARD con scaler (para FP16)
            scaler.scale(total_loss).backward()

            # STEP: actualizar parámetros
            scaler.step(optimizer)
            scaler.update()

            # Acumular losses para reporting
            train_ranking_loss_list.append(ranking_loss.item())
            train_identity_loss_list.append(identity_loss.item())
            train_total_loss_list.append(total_loss.item())

            tepoch.set_postfix({
                'rank': f"{np.mean(train_ranking_loss_list):.3f}",
                'id':   f"{np.mean(train_identity_loss_list):.3f}",
                'tot':  f"{np.mean(train_total_loss_list):.3f}",
            })

    scheduler.step()

    # Guardar checkpoint
    ckpt_path = os.path.join(
        config.model_save_path,
        f"TextReIDNet_epoch{current_epoch}.pth.tar"
    )
    torch.save({
        'epoch': current_epoch,
        'model_state_dict': model.state_dict(),
        'optimizer_state_dict': optimizer.state_dict(),
        'loss': np.mean(train_total_loss_list),
    }, ckpt_path)

    latest_path = os.path.join(
        config.model_save_path, "TextReIDNet_latest.pth.tar"
    )
    torch.save({
        'epoch': current_epoch,
        'model_state_dict': model.state_dict(),
        'optimizer_state_dict': optimizer.state_dict(),
        'loss': np.mean(train_total_loss_list),
    }, latest_path)
\end{minted}

\section{Estimación de tiempo}

Con los parámetros default (60 epochs, batch 8, FP16):

\begin{itemize}
  \item \textbf{RTX 2060}: 2--4 días.
  \item \textbf{RTX 3090}: 16--24 horas.
  \item \textbf{A100}: 8--12 horas.
\end{itemize}

\noindent\textbf{Nota}: el HF dataset tiene 179k muestras (vs 40k
del original). El entrenamiento es proporcionalmente más largo
que el paper.

\section{Caveats}

\begin{enumerate}
  \item \textbf{Sin early stopping}: el código entrena los 60
        epochs completos. \texttt{patience = 3} está en config
        pero no se usa.
  \item \textbf{Sin validación intermedia}: no evalúa sobre
        val/test durante el entrenamiento.
  \item \textbf{Sin accumulation de gradientes}: si OOM, hay
        que bajar \texttt{batch\_size}.
  \item \textbf{Checkpoint por epoch}: 61 archivos
        \texttt{.pth.tar} ($\sim$250 MB cada uno para 25M
        parámetros).
\end{enumerate}
